% =========================================================
% Proof Stub: Cauchy Product Theorem
% Source: volume-iii/analysis/sequences/notes/series/notes-series-rearrangements.tex
% =========================================================

% *** HARDER PROOF — attempt after foundational results settled ***

\clearpage
\phantomsection
\hypertarget{proof-RA-ABB-T-cauchy-product}{}

\subsubsection[Cauchy Product Theorem]{Proof --- Cauchy Product Theorem}

\bigskip

\noindent
\textbf{Source.} See \texttt{notes-series-rearrangements.tex}.

\vspace{0.75em}

\noindent
\textbf{Goal.} If $\sum a_n = A$ and $\sum b_n = B$ both converge absolutely, then the Cauchy product $\sum c_n$ (where $c_n = \sum_{k=0}^n a_k b_{n-k}$) converges absolutely to $AB$.

\vspace{0.75em}

\noindent
\textbf{Logical form.}
\[
  \sum |a_n| < \infty,\; \sum |b_n| < \infty \Rightarrow \sum c_n = AB.
\]

\vspace{0.75em}

\noindent
\textbf{Proof strategy.} Write $c_n$-partial sums as a truncated double sum over $\{(i,j) : i+j \le n\}$. The absolute convergence lets us rearrange the double series; the limit equals the product of the two individual sums.

\vspace{0.75em}

\noindent
\textbf{Proof.}
\begin{proof}
\Claim If $\sum a_n = A$ and $\sum b_n = B$ both converge absolutely, then the Cauchy product $\sum c_n$ (where $c_n = \sum_{k=0}^n a_k b_{n-k}$) converges absolutely to $AB$.

\Given 

\Goal 

\Strategy Write $c_n$-partial sums as a truncated double sum over $\{(i,j) : i+j \le n\}$. The absolute convergence lets us rearrange the double series; the limit equals the product of the two individual sums.

\medskip

% --- let A_n = ∑_{k=0}^n a_k, B_n = ∑_{k=0}^n b_k, C_n = ∑_{k=0}^n c_k ---
% --- C_n = ∑_{i+j≤n} a_i b_j ---
% --- A_n B_n = ∑_{i≤n, j≤n} a_i b_j = C_n + E_n (where E_n involves cross terms) ---
% --- show |E_n| ≤ (∑_{k>n/2}|a_k|)(∑|b_k|) + ... → 0 ---
% --- A_n B_n → AB and E_n → 0, so C_n → AB ---

\end{proof}

\begin{remark*}[Proof shape]
Bound the error between $A_n B_n$ and $C_n$ using absolute convergence; both tail sums go to 0.
\end{remark*}

\begin{remark*}[Dependencies]
\begin{itemize}
  \item Absolute convergence of both series.
  \item Algebra of limits.
  \item Triangle inequality for error estimation.
\end{itemize}
\end{remark*}
